\section{Construcción del panel de precios}
\label{anexo:datos}

Este anexo documenta las decisiones de construcción resumidas en la subsección
de Datos. Se detallan aquí porque condicionan lo que el diseño puede
identificar, pero su discusión completa interrumpiría la lectura del cuerpo.

\subsection{Identidad de la estación física}

La unidad de análisis es la estación física y no el identificador
administrativo de la CNE, porque un mismo local puede aparecer en el registro
bajo más de un código. Se detectaron dos situaciones, de naturaleza distinta.

\paragraph{Sucesión.} Un identificador nuevo comienza a reportar en el mismo
punto donde otro dejó de hacerlo. El criterio aplicado es que la primera fecha
del identificador entrante sea igual o posterior a la última del saliente y que
ambos disten menos de cincuenta metros. Corresponde típicamente a un cambio de
operador del mismo local. Si no se fusionaran, el cambio de operador se
contaría como una salida seguida de una entrada, es decir, se fabricaría un
evento de tratamiento donde no hubo ninguna alteración de la estructura
competitiva local.

\paragraph{Duplicación.} Los archivos del régimen antiguo llevan, para 31
estaciones, un segundo identificador formado agregando el sufijo \texttt{a} al
código base. La evidencia de que se trata del mismo local es concluyente: en
los 31 casos el código base también está presente en el registro, en los 31
casos ambos identificadores coexisten desde el inicio de la serie —no se
suceden— y los 31 dejan de reportarse cuando la CNE migra al formato actual,
que incorpora una columna de modalidad de atención. Es decir, lo que el régimen
antiguo llevaba como dos registros separados, el nuevo lo consolida en un solo
registro con dos modalidades.

La duplicación es más dañina que la sucesión para este trabajo. Dos registros
del mismo local aparecen en el panel como dos estaciones distintas separadas
por una distancia nula, de modo que cada una cuenta a la otra como competidora
inmediata. Esto no solo infla el número de estaciones del mercado, sino que
contamina precisamente la variable —el número de competidores en un radio
dado— que el diseño usa para delimitar los mercados locales.

Ambas situaciones se resuelven encadenando los identificadores hasta una raíz
común, que define la estación física. Tras la fusión, el registro contiene
2.030 estaciones físicas.

\subsection{Frecuencia del panel}

El panel se construye a frecuencia semanal sobre la grilla de jueves del MEPCO
y de ahí se colapsa a frecuencia mensual, tomando la última semana de cada mes.

La elección de la semana responde a dos hechos. El primero es la cadencia
efectiva de reporte: la mediana es de cincuenta observaciones por estación,
combustible y año, es decir, aproximadamente una por semana. El segundo es que
el precio mayorista con MEPCO se fija semanalmente y rige desde el jueves, de
modo que la semana es la unidad natural del costo.

Construir el panel a frecuencia diaria habría sido posible, pero cerca del
$85\%$ de sus celdas correspondería a un valor arrastrado desde el último
reporte, sin información adicional, y con observaciones perfectamente
correlacionadas dentro de cada semana. En la grilla semanal, en cambio, el
$81{,}1\%$ de las celdas corresponde a un precio efectivamente informado esa
semana.

\subsection{Arrastre y ventanas de actividad}

Dentro de cada estación y combustible el precio se arrastra hacia adelante
hasta un máximo de trece semanas, aproximadamente tres meses. El arrastre
refleja el hecho de que un precio no informado sigue vigente: la obligación
legal es informar cada variación, de modo que la ausencia de reporte significa
que el precio no cambió.

El tope cumple la función opuesta. Un vacío superior a trece semanas deja a la
estación fuera del panel durante esas semanas, de manera que un precio no se
arrastra a través de un período en que la estación no operaba o no reportaba.
Sin este tope, una estación cerrada durante un año seguiría figurando en el
panel con su último precio conocido, y sería contada como competidora activa de
sus vecinas.

\subsection{Márgenes}

El margen bruto se define como el precio de venta menos el precio mayorista con
MEPCO vigente esa semana, conforme a la \Cref{prop:margen}. Al ser el MEPCO una
referencia nacional fijada por regla, el costo se adjudica por combustible y
semana y no por estación.

La variable existe para el $82{,}4\%$ de las observaciones del panel. El resto
se reparte entre dos causas conocidas: la gasolina de 95 octanos, para la que no
se publica referencia MEPCO, y los meses anteriores a agosto de 2014, cuando el
mecanismo todavía no operaba. Ninguna de las dos introduce selección, ya que
ambas son ausencias por combustible y por período, no por estación.

Las series descriptivas de precio y margen del cuerpo excluyen los meses que el
panel no abarca de extremo a extremo. En la práctica esto descarta julio de
2026: el registro termina el día 12 de ese mes y el promedio queda construido
sobre dos semanas que además caen en medio de un traspaso inacabado, ya que el
precio mayorista había caído 95 pesos por litro el 18 de junio y el precio
minorista continuaba ajustándose. El margen promedio de ese mes es un artefacto
de truncamiento, no un resultado de mercado.

\subsection{Restricciones sobre la definición de entrada}

Una entrada se define como el primer mes en que una estación física aparece en
el registro. Sobre esa definición operan tres restricciones, que son las que
determinan qué eventos pueden usarse como tratamiento.

\paragraph{Censura por la izquierda.} Las estaciones presentes desde el inicio
de la serie no constituyen entradas: 1.548 de las 2.030 estaciones ya reportaban
en 2012 y su fecha de apertura es desconocida. Tratarlas como entrantes en enero
de 2012 fecharía como evento lo que es apenas el comienzo de la observación.

\paragraph{La cohorte de 2013.} Ese año se registran 79 entradas, contra veinte
a cincuenta en un año normal, y 43 de ellas corresponden a estaciones de bandera
blanca, cifra que no se repite en ningún otro año de la serie. Es el patrón de
una incorporación tardía al sistema de reporte —estaciones pequeñas que se
suman a Bencina en Línea cuando la plataforma ya está en marcha— y no el de una
ola de aperturas. Estos eventos se conservan para efectos de contaminación y de
corte de las ventanas, de modo que un incumbente expuesto a una de ellas no se
usa como si nada le hubiera ocurrido, pero nunca definen una cohorte de
tratamiento.

\paragraph{Solo la primera entrada.} De cada incumbente se considera únicamente
la primera entrada cercana, y las observaciones posteriores a una segunda se
descartan. El estimando corresponde así al efecto de una primera entrada y no a
una mezcla de efectos de primera y segunda, que no tienen por qué ser iguales:
el paso de un competidor a dos altera la estructura del mercado local de manera
distinta que el paso de dos a tres.

\subsection{Salidas y el quiebre de régimen de reporte}

El diseño empírico se construye exclusivamente sobre entradas, pero conviene
documentar por qué no se utilizan las salidas.

El cambio entre los archivos del régimen antiguo y los del actual deja una
cicatriz nítida: cien estaciones dejan de ser informadas en diciembre de 2022,
contra una o dos al año en todo el período posterior. Una parte de ese salto
corresponde a las duplicaciones ya descritas. El resto son, con alta
probabilidad, estaciones que habían dejado de actualizar sus precios en alguna
fecha anterior no observada y que sobrevivían en los archivos anuales con su
último precio conocido hasta que la migración las depuró.

La consecuencia es que la fecha de salida de esas estaciones no es
interpretable: lo único que se observa es el momento en que el sistema dejó de
listarlas, que no coincide con el momento en que dejaron de operar. Por eso las
tablas descriptivas reportan diciembre de 2022 en una categoría separada, en
lugar de diluirlo dentro de un total anual que sugeriría una ola de cierres que
no ocurrió.

Esta limitación no se traslada al diseño. La primera aparición de una estación
en el registro sí es una fecha interpretable, porque una estación que no existe
no puede informar precios. Un ejercicio de verificación confirma que las
entradas no arrastran el problema: de las 482 entradas identificadas, ninguna
tiene una estación que haya dejado de reportar dentro de los cincuenta metros
en los dieciocho meses previos, y solo siete la tienen dentro de doscientos
metros.

\subsection{Reemplazo de estaciones salientes}

Una precisión sobre cómo debe leerse la información de reemplazo. Dado que la
fusión de identificadores descrita más arriba ya une a una estación con su
sucesora en el mismo sitio, un reemplazo en el mismo local no puede aparecer en
los datos como una salida: aparece como continuidad de la misma estación
física. La pregunta que sí tiene contenido empírico es, entonces, si el mercado
local se repuso, y se responde verificando si abrió alguna estación a menos de
quinientos metros dentro de los veinticuatro meses siguientes al cierre. La
respuesta es que casi nunca ocurre, lo que resulta consistente con las barreras
a la entrada que documenta la \textcite{FNE2021informe}.
